% from ega3-1-fr.tex
\oldpage[0\textsubscript{III}]{49}
\[
\label{0.11.10.1.2-fr}
\tag{11.10.1.2}
\chi(M^\bullet)=\sum_{i=0}^{n}(-1)^i\operatorname{long}(M^i).
\]
\end{env}

\begin{proposition}[11.10.2]
\label{0.11.10.2-fr}
\emph{Pour tout complexe fini $M^\bullet$ de $A$-modules à gauche de longueur
finie, on a
\[
\chi(M^\bullet)=\chi(H^\bullet(M^\bullet))
\]
($H^\bullet(M^\bullet)$ étant considéré comme un complexe pour la dérivation
triviale). En particulier, si la suite
(\hyperref[0.11.10.1.1-fr]{11.10.1.1}) est exacte, on a
$\chi(M^\bullet)=0$.}
\end{proposition}

\begin{proof}
Posons pour abréger $B^i=B^i(M^\bullet)$, $Z^i=Z^i(M^\bullet)$,
$H^i=H^i(M^\bullet)=Z^i/B^i$; les $B^i$, $Z^i$, $H^i$ sont de longueur
finie. Des suites exactes
\[
0\longrightarrow B^i\longrightarrow Z^i\longrightarrow H^i
\longrightarrow0,
\qquad
0\longrightarrow Z^i\longrightarrow M^i\longrightarrow B^{i+1}
\longrightarrow0
\]
on tire les relations
\[
\operatorname{long}(Z^i)=\operatorname{long}(H^i)+
\operatorname{long}(B^i),
\qquad
\operatorname{long}(M^i)=\operatorname{long}(Z^i)+
\operatorname{long}(B^{i+1}),
\]
d'où
\[
\operatorname{long}(M^i)-\operatorname{long}(H^i)=
\operatorname{long}(B^{i+1})+\operatorname{long}(B^i).
\]
Multiplions cette relation par $(-1)^i$ et faisons la somme des relations
obtenues pour $0\leq i\leq n$, en notant que
$B^0=B^{n+1}=0$, il vient l'égalité cherchée.
\end{proof}

\begin{corollary}[11.10.3]
\label{0.11.10.3-fr}
Soit $E=(E_r^{pq})$ une suite spectrale dans la catégorie des modules sur un
anneau $A$. On suppose que les $E_2^{pq}$ sont des $A$-modules de longueur
finie et qu'il n'y a qu'un nombre fini de couples $(p,q)$ tels que
$E_2^{pq}\neq0$. Alors les caractéristiques d'Euler-Poincaré
$\chi(E_r^\bullet)$ de tous les complexes
$E_r^\bullet=(E_r^{(n)})_{n\in\mathbb Z}$
(\hyperref[0.11.1.1-fr]{11.1.1}) sont toutes égales. Si, en outre, la suite
$E$ est faiblement convergente et si on pose
\[
E_\infty^{(n)}=\bigoplus_{p+q=n}E_\infty^{pq}\quad(n\in\mathbb Z),
\]
on a aussi $\chi(E_\infty^\bullet)=\chi(E_2^\bullet)$,
$E_\infty^\bullet=(E_\infty^{(n)})_{n\in\mathbb Z}$ étant considéré comme
un complexe à dérivation triviale.
\end{corollary}

\begin{proof}
Notons d'abord que si $E_2^{pq}=0$, on a $E_r^{pq}=0$ pour
$2\leq r<+\infty$, donc tous les complexes $E_r^\bullet$ sont finis et
formés de $A$-modules de longueur finie; la relation
$\chi(E_r^\bullet)=\chi(E_{r+1}^\bullet)$ pour tout $r$ fini résulte donc de
(\hyperref[0.11.10.2-fr]{11.10.2}) et de l'isomorphie entre
$H^\bullet(E_r^\bullet)$ et $E_{r+1}^\bullet$ (en tant que complexes à
dérivation triviale). L'hypothèse que les $E_r^{pq}$ sont de longueur finie
entraîne que pour tout couple $(p,q)$, les suites
$(B_k(E_2^{pq}))_{k\geq2}$ et $(Z_k(E_2^{pq}))_{k\geq2}$ sont stationnaires;
l'hypothèse que $E$ est faiblement convergente et que $E_2^{pq}=0$, sauf
pour un nombre fini de couples $(p,q)$, entraîne donc qu'il existe un entier
$r\geq2$ tel que $E_\infty^{pq}=E_r^{pq}$ pour tout couple $(p,q)$; d'où
l'assertion relative à $\chi(E_\infty^\bullet)$.
\end{proof}

\section{Compléments sur la cohomologie des faisceaux}
\label{section:0.12-fr}

\subsection{Cohomologie des faisceaux de modules sur les espaces annelés}
\label{subsection:0.12.1-fr}

\begin{env}[12.1.1]
\label{0.12.1.1-fr}
Soit $(X,\mathcal O_X)$ un espace annelé. Rappelons que pour tout
$\mathcal O_X$-Module $\mathcal F$, on définit la cohomologie
$H^\bullet(X,\mathcal F)$, qui est un foncteur cohomologique universel
(T, 2.2) de la catégorie $C(X)$ des $\mathcal O_X$-Modules dans la catégorie
des groupes abéliens; c'est le foncteur dérivé du foncteur exact à gauche
$\mathcal F\to\Gamma(X,\mathcal F)$. Le foncteur $H^\bullet(X,\mathcal F)$
est isomorphe
